\chapter[23]{Heat Kernel for Static Metrics}
\label{chow-iii-ch23-heat-kernel-static-metrics}

\section{Fundamental solutions and the parametrix}

\begin{definition}[Fundamental solution to the heat equation]
\label{def:chow-iii-ch23-fundamental-solution}
\source{chow-et-al-ricci-flow-part-iii:page-0236}
\dcref{ch23:23.1}
\group{Heat kernels}

Let $(\mathcal M^n,g)$ be complete and put
\(\mathbb R^2_>:=\{(t,u):t>u\}\).  A function
\(h:\mathcal M\times\mathcal M\times\mathbb R^2_>\to\mathbb R\) is a
fundamental solution if it is continuous, $C^2$ in the space variables and
$C^1$ in the time variables, satisfies
\[
(\partial_t-\Delta_x)h(\cdot,\cdot;y,u)=0,
\qquad (\partial_u+\Delta_y)h(x,t;\cdot,\cdot)=0,
\]
and has the distributional limits
\[
\lim_{t\searrow u}h(\cdot,t;y,u)=\delta_y,
\qquad \lim_{u\nearrow t}h(x,t;\cdot,u)=\delta_x.
\]
Equivalently, for every compactly supported continuous $f$,
\[
\lim_{t\searrow u}\int h(x,t;y,u)f(x)\,d\mu(x)=f(y),\qquad
\lim_{u\nearrow t}\int h(x,t;y,u)f(y)\,d\mu(y)=f(x).
\]
\end{definition}

\begin{remark}[Completeness]
\label{rem:chow-iii-ch23-completeness}
\source{chow-et-al-ricci-flow-part-iii:page-0236}
\dcref{ch23:23.2}
\group{Heat kernels}

Completeness is not needed for the definition, but is assumed in the
applications below.
\end{remark}

\begin{theorem}[Existence of the heat kernel on a closed manifold]
\label{thm:chow-iii-ch23-closed-heat-kernel}
\source{chow-et-al-ricci-flow-part-iii:page-0237}
\dcref{ch23:23.3}
\group{Heat kernels}

If $(\mathcal M^n,g)$ is closed, there is a unique positive smooth symmetric
fundamental solution $H(x,t;y,u)$ and
\[
\int_{\mathcal M}H(x,t;y,u)\,d\mu(x)=1.
\]
\end{theorem}
\begin{proof}
\sketch The construction in Theorem~\ref{thm:chow-iii-ch23-parametrix-heat-kernel}
produces a smooth fundamental solution.  Uniqueness, positivity, and symmetry
follow from the standard energy and maximum-principle arguments for the heat
equation on a closed manifold; integrating the equation in $x$ gives the
constant mass, whose limiting value is one.
\end{proof}

\begin{remark}[Properties treated elsewhere]
\label{rem:chow-iii-ch23-properties-elsewhere}
\source{chow-et-al-ricci-flow-part-iii:page-0237}
\dcref{ch23:23.4}
\group{Heat kernels}

The parametrix argument below proves existence and smoothness; uniqueness,
positivity, and symmetry also follow from the earlier heat-kernel theory.
\end{remark}

\begin{definition}[Transplanted Euclidean heat kernel]
\label{def:chow-iii-ch23-transplanted-kernel}
\source{chow-et-al-ricci-flow-part-iii:page-0238}
\group{Parametrix}

For $d(x,y)<\operatorname{inj}(g)$ set
\[
E(x,t;y,u):=(4\pi(t-u))^{-n/2}
 \exp\!\left(-\frac{d(x,y)^2}{4(t-u)}\right).
\]
It is smooth off the cut locus and is the Euclidean heat kernel when
$\mathcal M=\mathbb R^n$.
\end{definition}

\begin{remark}[Choice of the finite approximation]
\label{rem:chow-iii-ch23-finite-approximation}
\source{chow-et-al-ricci-flow-part-iii:page-0238}
\group{Parametrix}

For $N>n/2$ we seek
\[
H_N(x,t;y,u)=E(x,t;y,u)\sum_{k=0}^N\phi_k(x,y)(t-u)^k
\]
on the injectivity neighborhood.  The last term tends uniformly to zero as
$t-u\to0$ because $N>n/2$.
\end{remark}

\begin{definition}[Normal-coordinate Jacobian]
\label{def:chow-iii-ch23-normal-jacobian}
\source{chow-et-al-ricci-flow-part-iii:page-0239}
\group{Parametrix}

For $r(x)=d(x,y)$ and geodesic spherical coordinates about $y$, define
\[
\alpha(x,y):=\frac{\sqrt{\det g^S(x)}}{r(x)^{n-1}}.
\]
In normal coordinates this is $\sqrt{\det(g_{ij})(x)}$ and extends smoothly
to the diagonal with value one.
\end{definition}

\begin{remark}[Gauss lemma and the radial Laplacian]
\label{rem:chow-iii-ch23-gauss-lemma}
\source{chow-et-al-ricci-flow-part-iii:page-0240}
\dcref{ch23:23.6}
\group{Parametrix}

The Gauss lemma gives $x^i=\sum_jx^jg_{ji}$ and, for a radial function $f(r)$,
\[
\Delta_xf=f''+\left(\frac{n-1}{r}+\partial_r\log\alpha\right)f'.
\]
Consequently
\[
(\Delta_x-\partial_t)E=-\frac r{2(t-u)}(\partial_r\log\alpha)E.
\]
\end{remark}

\begin{lemma}[Recursive ODEs for the coefficients]
\label{lem:chow-iii-ch23-recursive-odes}
\source{chow-et-al-ricci-flow-part-iii:page-0241,chow-et-al-ricci-flow-part-iii:page-0242}
\uses{def:chow-iii-ch23-normal-jacobian}
\dcref{ch23:23.7}
\group{Parametrix}


Let $\phi_0(y,y)=1$.  The identity
\[
(\Delta_x-\partial_t)H_N=E(\Delta_x\phi_N)(t-u)^N
\]
is equivalent along radial geodesics to
\[
\partial_r\phi_0+\tfrac12(\partial_r\log\alpha)\phi_0=0,
\]
and, for $1\le k\le N$,
\[
r\partial_r\phi_k+\left(\tfrac r2\partial_r\log\alpha+k\right)\phi_k
=\Delta_x\phi_{k-1}.
\]
The finite solutions are smooth on the injectivity neighborhood and satisfy
\(\phi_0=\alpha^{-1/2}\).
\end{lemma}
\begin{proof}
\sketch Apply the radial Laplacian formula to $E\sum_k\phi_k(t-u)^k$.  Equating each
power of $t-u$ gives the displayed ODEs.  Multiplication of the $k$th equation
by $r^{k-1}\alpha^{1/2}$ gives an exact derivative; integration from the
diagonal, with finiteness there, yields a smooth solution recursively.
\end{proof}

\begin{remark}[Independence of $N$]
\label{rem:chow-iii-ch23-coefficients-independent}
\source{chow-et-al-ricci-flow-part-iii:page-0242}
\dcref{ch23:23.8}
\group{Parametrix}

Each $\phi_k$ is independent of the choice of $N\ge k$.
\end{remark}

\begin{lemma}[Recursive formula for the coefficients]
\label{lem:chow-iii-ch23-recursive-formula}
\source{chow-et-al-ricci-flow-part-iii:page-0242}
\uses{lem:chow-iii-ch23-recursive-odes}
\dcref{ch23:23.9}
\group{Parametrix}


If the coefficients are finite on the diagonal, then
\[
\phi_k(x,y)=\alpha(x,y)^{-1/2}\int_0^1\rho^{k-1}
 \bigl(\alpha^{1/2}\Delta_x\phi_{k-1}\bigr)
 \bigl(\exp_y(\rho r(x)V),y\bigr)\,d\rho,
\]
where $V$ is the initial unit tangent of the minimizing geodesic from $y$ to
$x$.
\end{lemma}
\begin{proof}
\sketch Integrate the exact derivative furnished by
Lemma~\ref{lem:chow-iii-ch23-recursive-odes} and substitute
$r=\rho r(x)$.
\end{proof}

\begin{lemma}[Covariant derivatives of the heat error]
\label{lem:chow-iii-ch23-heat-error-derivatives}
\source{chow-et-al-ricci-flow-part-iii:page-0243,chow-et-al-ricci-flow-part-iii:page-0244,chow-et-al-ricci-flow-part-iii:page-0245}
\uses{lem:chow-iii-ch23-recursive-formula}
\dcref{ch23:23.10}
\group{Parametrix}


For every $k,\ell\ge0$ and $T>0$,
\[
|\partial_t^\ell\nabla_x^k(\Box_xH_N)|
\le C(t-u)^{N-n/2-k-2\ell}
 e^{-d(x,y)^2/4(t-u)}.
\]
If $N>n/2+k+2\ell$, then the left side equals
\[
(4\pi)^{-n/2}(t-u)^{N-n/2-k-2\ell}e^{-d^2/4(t-u)}
 \widetilde F_{k,\ell}(x,y,t-u)
\]
for a smooth covariant tensor $\widetilde F_{k,\ell}$.
\end{lemma}
\begin{proof}
\sketch Differentiate $(\Box_xH_N)=E(\Delta_x\phi_N)(t-u)^N$.  Each spatial derivative
costs at most one power of $(t-u)^{-1}$ and each time derivative at most two.
On a finite atlas, boundedness of the coefficients and their derivatives gives
the asserted uniform estimate and the smooth remainder tensor.
\end{proof}

\begin{definition}[Parametrix for the heat operator]
\label{def:chow-iii-ch23-parametrix}
\source{chow-et-al-ricci-flow-part-iii:page-0245}
\dcref{ch23:23.11}
\group{Parametrix}

A smooth $P:\mathcal M\times\mathcal M\times\mathbb R^2_>\to\mathbb R$ is a
parametrix if $(\Delta_x-\partial_t)P$ and $(\Delta_y+\partial_u)P$ extend
continuously to $t=u$, and both distributional initial limits equal the
corresponding Dirac mass.
\end{definition}

\begin{proposition}[Existence of a parametrix]
\label{prop:chow-iii-ch23-parametrix-existence}
\source{chow-et-al-ricci-flow-part-iii:page-0246,chow-et-al-ricci-flow-part-iii:page-0247}
\uses{def:chow-iii-ch23-parametrix, lem:chow-iii-ch23-heat-error-derivatives}
\dcref{ch23:23.12}
\group{Parametrix}


Choose a smooth radial $\eta$ with $\eta=1$ on $[0,\operatorname{inj}(g)/4]$,
$\eta=0$ on $[\operatorname{inj}(g)/2,\infty)$, bounded derivatives, and set
\[
P_N(x,t;y,u):=\eta(d(x,y))H_N(x,t;y,u).
\]
For $N>n/2$, $P_N$ is a parametrix.
\end{proposition}
\begin{proof}
\sketch The support of $\eta$ lies inside the injectivity neighborhood, so $P_N$ is
smooth.  The product rule gives
\[
\Box_xP_N=\eta\Box_xH_N+(\Delta_x\eta)H_N+2\langle\nabla\eta,\nabla H_N\rangle.
\]
Near the diagonal the first term tends to zero by Lemma~\ref{lem:chow-iii-ch23-heat-error-derivatives}; away from it every term is exponentially small.  The same argument in the $y$ variable proves continuous extension.  The Gaussian approximate-identity limit follows by normal coordinates and the fact that $\eta=1$ near the diagonal; symmetry gives the second limit.
\end{proof}

\begin{example}[Convolution and sum]
\label{ex:chow-iii-ch23-convolution-sum}
\source{chow-et-al-ricci-flow-part-iii:page-0248}
\dcref{ch23:23.13}
\group{Parametrix}

For $\Box_xP_N$ one has
\[
(\Box_xP_N)*\sum_{k\ge1}(\Box_xP_N)^{*k}
=\sum_{k\ge2}(\Box_xP_N)^{*k}.
\]
\end{example}
\begin{proof}
\sketch Expand the convolution, use bilinearity and associativity, and pass to the
limit using the locally uniform absolute convergence established below.
\end{proof}

\begin{lemma}[Derivatives of the heat error of the parametrix]
\label{lem:chow-iii-ch23-parametrix-error-derivatives}
\source{chow-et-al-ricci-flow-part-iii:page-0248,chow-et-al-ricci-flow-part-iii:page-0249}
\uses{prop:chow-iii-ch23-parametrix-existence}
\dcref{ch23:23.14}
\group{Parametrix}


For every $k,\ell\ge0$,
\[
\partial_t^\ell\nabla_x^k(\Box_xP_N)
=(t-u)^{N-n/2-k-2\ell}e^{-d(x,y)^2/5(t-u)}G_{k,\ell}(x,y,t-u),
\]
where $G_{k,\ell}$ is smooth and supported in the injectivity neighborhood.
\end{lemma}
\begin{proof}
\sketch Inside the region where $\eta=1$ use Lemma~\ref{lem:chow-iii-ch23-heat-error-derivatives}; on the annulus where derivatives of $\eta$ occur the distance is bounded below, so the Gaussian dominates every negative power of $t-u$.  The exterior is identically zero.
\end{proof}

\section{Convolution construction}

\begin{definition}[Space-time convolution]
\label{def:chow-iii-ch23-space-time-convolution}
\source{chow-et-al-ricci-flow-part-iii:page-0250}
\dcref{ch23:eq-23.46,eq-23.47,eq-23.48,eq-23.49}
\group{Convolution}

For continuous $F,G$ define
\[
(F*G)(x,y,t):=\int_0^t\int_{\mathcal M}F(x,z,s)G(z,y,t-s)\,d\mu(z)\,ds.
\]
The operation is bilinear and associative; write $F^{*k}$ for the $k$-fold
convolution.
\end{definition}
\begin{proof}
\sketch Associativity follows by expanding both sides as a triple integral and applying
Fubini on the compact manifold, followed by the change of variables
$(r,s)\leftrightarrow(s-r,t-s)$.
\end{proof}

\begin{lemma}[Heat operator of a convolution with the parametrix]
\label{lem:chow-iii-ch23-parametrix-convolution-heat-operator}
\source{chow-et-al-ricci-flow-part-iii:page-0251}
\uses{def:chow-iii-ch23-space-time-convolution, prop:chow-iii-ch23-parametrix-existence}
\dcref{ch23:23.15}
\group{Convolution}


For continuous $G$,
\[
\Box_x(P_N*G)=(\Box_xP_N)*G-G.
\]
\end{lemma}
\begin{proof}
\sketch Differentiate the defining integral in $t$ and use the upper-endpoint term
from the approximate-identity property of $P_N$; differentiating in $x$ is
legitimate away from the endpoint and then follows by the estimates in
Lemma~\ref{lem:chow-iii-ch23-parametrix-error-derivatives}.
\end{proof}

\begin{theorem}[Existence by the parametrix convolution series]
\label{thm:chow-iii-ch23-parametrix-heat-kernel}
\source{chow-et-al-ricci-flow-part-iii:page-0252,chow-et-al-ricci-flow-part-iii:page-0253}
\uses{lem:chow-iii-ch23-parametrix-convolution-heat-operator, cor:chow-iii-ch23-convolution-series}
\dcref{ch23:23.16}
\group{Convolution}


For $N>n/2$, put
\[
G_N:=\sum_{k=1}^{\infty}(\Box_xP_N)^{*k},\qquad H:=P_N+P_N*G_N.
\]
Then $H$ is smooth, independent of $N$, satisfies $\Box_xH=0$, and is the
unique fundamental solution.
\end{theorem}
\begin{proof}
\sketch The convolution estimates below give locally uniform convergence of $G_N$ and
of all derivative series after increasing $N$.  Lemma~\ref{lem:chow-iii-ch23-parametrix-convolution-heat-operator} and the convolution-sum identity give
\(\Box_xH=\Box_xP_N+(\Box_xP_N)*G_N-G_N=0\).  The error term
\(P_N*G_N\) tends to zero against every continuous test function as $t\searrow u$,
so $H$ has the required initial limit.  The same argument in the second
variable gives the other limit.  Uniqueness of the fundamental solution on a
closed manifold makes the expression independent of $N$ and completes the
proof.
\end{proof}

\begin{theorem}[Short-time heat-kernel expansion]
\label{thm:chow-iii-ch23-heat-kernel-expansion}
\source{chow-et-al-ricci-flow-part-iii:page-0253,chow-et-al-ricci-flow-part-iii:page-0254}
\uses{thm:chow-iii-ch23-parametrix-heat-kernel, lem:chow-iii-ch23-convolution-estimates}
\dcref{ch23:23.17}
\group{Asymptotics}


There are smooth $u_j$ on $\mathcal M\times\mathcal M$, with $u_0(y,y)=1$, such that
\[
H(x,y,t)\sim(4\pi t)^{-n/2}e^{-d(x,y)^2/4t}\sum_{j=0}^{\infty}t^ju_j(x,y).
\]
More precisely, after truncation at $N$, the remainder is
\(O(t^{N+1-n/2})\) uniformly in $x,y$.
\end{theorem}
\begin{proof}
\sketch Take $u_j=\eta(d)\phi_j$ for $j\le N$ and subtract the finite parametrix
from the representation of $H$.  The convolution estimate in
Lemma~\ref{lem:chow-iii-ch23-convolution-estimates} bounds the difference by
\(C\int_0^t(t-s)^{N-n/2}\,ds=O(t^{N+1-n/2})\).
\end{proof}

\begin{example}[Differentiated expansion remainder]
\label{ex:chow-iii-ch23-expansion-derivatives}
\source{chow-et-al-ricci-flow-part-iii:page-0254}
\dcref{ch23:23.18}
\group{Asymptotics}

For $w_N$ the remainder above,
\[
|\partial_t^\ell\nabla_x^kw_N|=O(t^{N+1-n/2-k-2\ell})
\]
uniformly as $t\to0$.
\end{example}
\begin{proof}
\sketch Differentiate the convolution representation of $w_N$ and use the derivative
estimates for the parametrix error; the same time integral gains one power of
$t$.
\end{proof}

\begin{lemma}[Structure of the parametrix error]
\label{lem:chow-iii-ch23-coefficient-structure}
\source{chow-et-al-ricci-flow-part-iii:page-0254,chow-et-al-ricci-flow-part-iii:page-0255}
\dcref{ch23:23.19}
\group{Convolution}

For $N>n/2$,
\[
\Box_xP_N=t^N(4\pi t)^{-n/2}e^{-d(x,y)^2/5t}q_N(x,y,t),
\]
where $q_N$ is smooth on $\mathcal M\times\mathcal M\times[0,\infty)$.
\end{lemma}
\begin{proof}
\sketch Use the product formula for $P_N$ and factor the Gaussian with denominator five.
The terms containing derivatives of $\eta$ are supported where
$d(x,y)\ge\operatorname{inj}(g)/4$, hence are flat at $t=0$; all other terms
contain the factor $t^N$ by the defining equation for $H_N$.
\end{proof}

\begin{lemma}[Convolution estimates]
\label{lem:chow-iii-ch23-convolution-estimates}
\source{chow-et-al-ricci-flow-part-iii:page-0255,chow-et-al-ricci-flow-part-iii:page-0256}
\uses{lem:chow-iii-ch23-coefficient-structure}
\dcref{ch23:23.20}
\group{Convolution}


If $C=\sup|q_N|$ on $\mathcal M^2\times[0,T]$, then
\[
|(\Box_xP_N)^{*k}(x,y,t)|\le
\frac{C^k\operatorname{Vol}(\mathcal M)^{k-1}t^{k(N-n/2+1)-1}}
 {(k-1)!(N-n/2+1)^{k-1}}e^{-d(x,y)^2/5t}.
\]
\end{lemma}
\begin{proof}
\sketch The case $k=1$ is Lemma~\ref{lem:chow-iii-ch23-coefficient-structure}.  For
$k+1$, convolve the induction bound with the $k=1$ bound and use
\(d(x,z)^2/s+d(z,y)^2/(t-s)\ge d(x,y)^2/t\); the time integral is the beta
integral and gives the displayed factorial denominator.
\end{proof}

\begin{corollary}[Convergence of the convolution series]
\label{cor:chow-iii-ch23-convolution-series}
\source{chow-et-al-ricci-flow-part-iii:page-0256}
\uses{lem:chow-iii-ch23-convolution-estimates}
\dcref{ch23:23.21}
\group{Convolution}


The series $\sum_{k\ge1}e^{d(x,y)^2/5t}(\Box_xP_N)^{*k}$ converges absolutely and
uniformly on $\mathcal M^2\times(0,T]$; the unweighted series converges uniformly
up to $t=0$.
\end{corollary}
\begin{proof}
\sketch Sum the majorant in Lemma~\ref{lem:chow-iii-ch23-convolution-estimates}; it is
\(Ct^{N-n/2}\exp(C\operatorname{Vol}(\mathcal M)t^{N-n/2+1}/(N-n/2+1))\).
\end{proof}

\begin{example}[Commutation of convolution and the series]
\label{ex:chow-iii-ch23-convolution-series-commutes}
\source{chow-et-al-ricci-flow-part-iii:page-0257}
\uses{cor:chow-iii-ch23-convolution-series}
\dcref{ch23:23.22}
\group{Convolution}


The identity in Example~\ref{ex:chow-iii-ch23-convolution-sum} is valid because
the convergent series may be convolved term by term.
\end{example}

\begin{lemma}[Differentiability of the convolution series in coordinates]
\label{lem:chow-iii-ch23-convolution-series-derivatives}
\source{chow-et-al-ricci-flow-part-iii:page-0258}
\uses{lem:chow-iii-ch23-parametrix-error-derivatives, lem:chow-iii-ch23-convolution-estimates}
\dcref{ch23:23.23}
\group{Convolution}


On compact coordinate sets and $[0,T]$, if $N>n/2+2\ell+m$, then
\[
\sum_{k\ge1}\partial_t^\ell\partial_x^m(\Box_xP_N)^{*k}
\]
converges absolutely and uniformly, and equals
\(\partial_t^\ell\partial_x^mG_N\).
\end{lemma}
\begin{proof}
\sketch Differentiate the first factor of each convolution.  The derivative estimate
has time exponent $N-n/2-2\ell-m>-1$; convolving with the majorant in
Lemma~\ref{lem:chow-iii-ch23-convolution-estimates} yields a factorially summable
bound, uniformly on compact sets.
\end{proof}

\begin{lemma}[Covariant differentiability of the convolution series]
\label{lem:chow-iii-ch23-convolution-series-covariant}
\source{chow-et-al-ricci-flow-part-iii:page-0259}
\uses{lem:chow-iii-ch23-convolution-series-derivatives}
\dcref{ch23:23.24}
\group{Convolution}


Under the same hypotheses, the series of covariant $m$-tensors
\[
\sum_{k\ge1}\partial_t^\ell\nabla_x^m(\Box_xP_N)^{*k}
\]
converges uniformly on $\mathcal M^2\times[0,T]$ and equals
\(\partial_t^\ell\nabla_x^mG_N\).
\end{lemma}
\begin{proof}
\sketch In a finite atlas, covariant derivatives are linear combinations of coordinate
derivatives with smooth bounded coefficients, so Lemma~\ref{lem:chow-iii-ch23-convolution-series-derivatives} applies.
\end{proof}

\section{Differentiating convolutions}

\begin{lemma}[Continuity of integration against the parametrix]
\label{lem:chow-iii-ch23-parametrix-integration-continuity}
\source{chow-et-al-ricci-flow-part-iii:page-0259,chow-et-al-ricci-flow-part-iii:page-0260}
\uses{def:chow-iii-ch23-parametrix}
\dcref{ch23:23.25}
\group{Convolution differentiation}


If $f:\mathcal M\times[0,T]\to\mathbb R$ is continuous, then
\[
I_N(x,t,u)=\int P_N(x,z,t-u)f(z,u)\,d\mu(z)
\]
is continuous, converges uniformly to $f(x,u)$ as $t\searrow u$, and is smooth
in $(x,t)$ with derivatives obtained by differentiating under the integral.
\end{lemma}
\begin{proof}
\sketch Split the integral into a small normal ball and its complement.  The Gaussian
bound makes the complement exponentially small and the normal-coordinate change
of variables reduces the ball to the Euclidean approximate identity; dominated
convergence handles all derivatives.
\end{proof}

\begin{lemma}[First space derivatives of a parametrix convolution]
\label{lem:chow-iii-ch23-first-convolution-derivative}
\source{chow-et-al-ricci-flow-part-iii:page-0260,chow-et-al-ricci-flow-part-iii:page-0261,chow-et-al-ricci-flow-part-iii:page-0262,chow-et-al-ricci-flow-part-iii:page-0263}
\uses{lem:chow-iii-ch23-parametrix-integration-continuity}
\dcref{ch23:23.26}
\group{Convolution differentiation}


For continuous $G$, $P_N*G$ is $C^1$ in $x$ and
\[
\partial_{x^i}(P_N*G)(x,y,t)=\int_0^t\int_{\mathcal M}
\partial_{x^i}P_N(x,z,t-s)G(z,y,s)\,d\mu(z)\,ds.
\]
\end{lemma}
\begin{proof}
\sketch The first derivative is bounded by
\(C(t-s)^{-\alpha}d(x,z)^{2\alpha-n-1}\) for $1/2<\alpha<1$,
which is integrable in $z$ and $s$.  Difference quotients, split away from and
near $s=t$, then converge by uniform continuity and this integrable bound.
\end{proof}

\begin{lemma}[Second space derivatives of a parametrix convolution]
\label{lem:chow-iii-ch23-second-convolution-derivative}
\source{chow-et-al-ricci-flow-part-iii:page-0263,chow-et-al-ricci-flow-part-iii:page-0264,chow-et-al-ricci-flow-part-iii:page-0265,chow-et-al-ricci-flow-part-iii:page-0266,chow-et-al-ricci-flow-part-iii:page-0267,chow-et-al-ricci-flow-part-iii:page-0268,chow-et-al-ricci-flow-part-iii:page-0269,chow-et-al-ricci-flow-part-iii:page-0270}
\uses{lem:chow-iii-ch23-first-convolution-derivative}
\dcref{ch23:23.27}
\group{Convolution differentiation}


If $B(x,\operatorname{inj}(g)/2)$ lies in a coordinate chart, then $P_N*G$ is
$C^2$ at $x$ and
\[
\partial_{x^i}\partial_{x^j}(P_N*G)=\int_0^t\int_{\mathcal M}
\partial_{x^i}\partial_{x^j}P_N(x,z,t-s)G(z,y,s)\,d\mu(z)\,ds.
\]
\end{lemma}
\begin{proof}
\sketch The difference between $x$- and $z$-second derivatives of $P_N$ is bounded by
\(C(1+d^2/(t-s))(t-s)^{-n/2}e^{-d^2/4(t-s)}\), hence is integrable.  Integrate
the $z$-derivative term by parts; its first derivative satisfies the integrable
bound used in Lemma~\ref{lem:chow-iii-ch23-first-convolution-derivative}.  Applying
the same difference-quotient argument proves the formula.
\end{proof}

\begin{corollary}[Laplacian of a parametrix convolution]
\label{cor:chow-iii-ch23-laplacian-convolution}
\source{chow-et-al-ricci-flow-part-iii:page-0264}
\uses{lem:chow-iii-ch23-second-convolution-derivative}
\dcref{ch23:23.28}
\group{Convolution differentiation}


In geodesic coordinates centered at $x$,
\[
\Delta_x(P_N*G)(x,y,t)=\int_0^t\int_{\mathcal M}
\Delta_xP_N(x,z,t-s)G(z,y,s)\,d\mu(z)\,ds.
\]
\end{corollary}

\begin{lemma}[Time derivative of a parametrix convolution]
\label{lem:chow-iii-ch23-time-convolution-derivative}
\source{chow-et-al-ricci-flow-part-iii:page-0270,chow-et-al-ricci-flow-part-iii:page-0271}
\uses{lem:chow-iii-ch23-parametrix-integration-continuity}
\dcref{ch23:23.29}
\group{Convolution differentiation}


For continuous $G$,
\[
\partial_t(P_N*G)(x,y,t)=G(x,y,t)+\int_0^t\int_{\mathcal M}
\partial_tP_N(x,z,t-s)G(z,y,s)\,d\mu(z)\,ds.
\]
\end{lemma}
\begin{proof}
\sketch Write the convolution as $\int_0^tJ_N(s,t)\,ds$.  The upper endpoint tends to
$G(x,y,t)$ by the approximate-identity lemma, while the difference quotient of
$J_N$ converges under the integral by the same Gaussian domination.
\end{proof}

\section{Static-metric asymptotics}

\begin{lemma}[Expansion of the volume density]
\label{lem:chow-iii-ch23-volume-density-expansion}
\source{chow-et-al-ricci-flow-part-iii:page-0272,chow-et-al-ricci-flow-part-iii:page-0273}
\dcref{ch23:23.30}
\group{Asymptotics}

\[
\alpha(x,y)=1-\frac16R_{pq}(y)x^px^q-
\frac1{12}\nabla_rR_{pq}(y)x^px^qx^r+O(r^4),
\]
and hence
\[
\phi_0=\alpha^{-1/2}=1+\frac1{12}R_{pq}x^px^q+
\frac1{24}\nabla_rR_{pq}x^px^qx^r+O(r^4).
\]
\end{lemma}
\begin{proof}
\sketch Insert the normal-coordinate metric expansion into the determinant expansion of
Lemma~\ref{lem:chow-iii-ch23-determinant-expansion}, contract curvature indices
to obtain the Ricci terms, and take the square root.
\end{proof}

\begin{lemma}[Laplacian of squared distance]
\label{lem:chow-iii-ch23-distance-square-expansion}
\source{chow-et-al-ricci-flow-part-iii:page-0274}
\uses{lem:chow-iii-ch23-volume-density-expansion}
\dcref{ch23:23.31}
\group{Asymptotics}


\[
\Delta_xd(x,y)^2=2n-\frac23R_{pq}(y)x^px^q-
\frac12\nabla_rR_{pq}(y)x^px^qx^r+O(r^4).
\]
\end{lemma}
\begin{proof}
\sketch Use $\Delta r=(n-1)/r+\partial_r\log\alpha$ and
\(\Delta r^2=2+2r\Delta r\), then differentiate the expansion for $\alpha$.
\end{proof}

\begin{lemma}[First derivatives of the leading coefficient]
\label{lem:chow-iii-ch23-leading-coefficient-derivatives}
\source{chow-et-al-ricci-flow-part-iii:page-0274,chow-et-al-ricci-flow-part-iii:page-0275}
\uses{lem:chow-iii-ch23-volume-density-expansion}
\dcref{ch23:23.32}
\group{Asymptotics}


\[
|\nabla_x\phi_0|^2=O(r^2),\qquad
\Delta_x\phi_0=\frac16R(y)+\frac16(\nabla_rR)(y)x^r+O(r^2).
\]
\end{lemma}
\begin{proof}
\sketch Differentiate the normal-coordinate expansion for $\alpha^{1/2}$ and use
\(\Delta f=\alpha^{-1}\partial_i(\alpha g^{ij}\partial_jf)\), together with
the contracted Bianchi identity.
\end{proof}

\begin{lemma}[First-order expansion of the next coefficient]
\label{lem:chow-iii-ch23-phi-one-expansion}
\source{chow-et-al-ricci-flow-part-iii:page-0276}
\uses{lem:chow-iii-ch23-recursive-formula, lem:chow-iii-ch23-leading-coefficient-derivatives}
\dcref{ch23:23.33}
\group{Asymptotics}


\[
\phi_1(x,y)=\frac16R(y)+\frac1{12}(\nabla_rR)(y)x^r+O(r^2).
\]
\end{lemma}
\begin{proof}
\sketch Insert the expansion of $\Delta\phi_0$ into the integral formula for $\phi_1$;
the substitution $x\mapsto\exp_y(\rho rV)$ replaces $x^r$ by $\rho x^r$ and
integrating $\rho$ from zero to one gives the coefficient $1/12$.
\end{proof}

\begin{remark}[The second diagonal coefficient]
\label{rem:chow-iii-ch23-phi-two-diagonal}
\source{chow-et-al-ricci-flow-part-iii:page-0276}
\uses{lem:chow-iii-ch23-recursive-formula}
\dcref{ch23:23.34}
\group{Asymptotics}


\[
\phi_2(y,y)=\frac1{30}\Delta R+\frac1{72}R^2-
\frac1{180}|\operatorname{Rc}|^2+\frac1{180}|\operatorname{Rm}|^2.
\]
\end{remark}

\begin{remark}[Static entropy asymptotics]
\label{rem:chow-iii-ch23-static-entropy}
\source{chow-et-al-ricci-flow-part-iii:page-0277,chow-et-al-ricci-flow-part-iii:page-0280}
\dcref{ch23:23.35}
\group{Asymptotics}

For $H_N=(4\pi t)^{-n/2}e^{-f}$, differentiating
\(\log H_N=-\frac n2\log(4\pi t)-d^2/4t+\log\sum_k\phi_kt^k\) and using the preceding expansions gives, on the diagonal,
\[
\left(\partial_t+\Delta_x\right)\!\left(\log H_N+\tfrac n2\log(4\pi t)\right)
 +\frac{\log H_N+\frac n2\log(4\pi t)}t
= -\frac n{2t}+\frac12R+O(t).
\]
\end{remark}

\begin{example}[Gaussian moment estimates]
\label{ex:chow-iii-ch23-gaussian-moments}
\source{chow-et-al-ricci-flow-part-iii:page-0278,chow-et-al-ricci-flow-part-iii:page-0279}
\dcref{ch23:eq-23.132,eq-23.133,eq-23.134}
\group{Asymptotics}

For the Euclidean Gaussian, odd errors of order $|x|$ and $|x|^3/t$ integrate to
$O(t^{1/2})$, and for symmetric $A$,
\[
\int_{\mathbb R^n}\frac{A_{ij}x^ix^j}{4t}(4\pi t)^{-n/2}e^{-|x|^2/4t}\,dx
=\frac{\operatorname{tr}A}{2}.
\]
The corresponding manifold estimates are uniform in the center $y$.
\end{example}
\begin{proof}
\sketch Scale $x=\sqrt{4t}\,\widetilde x$ and use rotational invariance; diagonalizing
$A$ reduces the last identity to the one-dimensional Gaussian second moment.
The manifold statements follow by normal coordinates and the heat-kernel
expansion.
\end{proof}

\section{Supplementary tools}

\begin{lemma}[Expansion for the determinant]
\label{lem:chow-iii-ch23-determinant-expansion}
\source{chow-et-al-ricci-flow-part-iii:page-0280,chow-et-al-ricci-flow-part-iii:page-0281,chow-et-al-ricci-flow-part-iii:page-0282}
\dcref{ch23:23.37}
\group{Supplementary analysis}

If
\[
M(s)=I+sA+s^2B+s^3C+s^4D+O(s^5),
\]
then
\[
\begin{aligned}
\det M={}&1+s\operatorname{tr}A+s^2\left(\operatorname{tr}B+	frac12(\operatorname{tr}A)^2-	frac12\operatorname{tr}(A^2)\right)\\
&+s^3\left(\operatorname{tr}A(\tfrac16(\operatorname{tr}A)^2-	frac12\operatorname{tr}(A^2)+\operatorname{tr}B)+\tfrac13\operatorname{tr}(A^3)-\operatorname{tr}(AB)+\operatorname{tr}C\right)+O(s^4).
\end{aligned}
\]
If $A=0$, then
\[
\det M=1+s^2\operatorname{tr}B+s^3\operatorname{tr}C+s^4\left(\tfrac12(\operatorname{tr}B)^2-	frac12\operatorname{tr}(B^2)+\operatorname{tr}D\right)+O(s^5).
\]
\end{lemma}
\begin{proof}
\sketch Differentiate $\det M$ using
\( (\det M)'=\det M\operatorname{tr}(M^{-1}M')\) through fourth order and
compare coefficients in Taylor's formula.  Setting $A=0$ gives the second
display.
\end{proof}

\begin{lemma}[Fubini theorem on noncompact space-time]
\label{lem:chow-iii-ch23-fubini-noncompact}
\source{chow-et-al-ricci-flow-part-iii:page-0282,chow-et-al-ricci-flow-part-iii:page-0283}
\dcref{ch23:23.38}
\group{Supplementary analysis}

If $f\in C^1(\mathcal M\times[\alpha,\omega])$ on a noncompact manifold and
\(\int_{\mathcal M}f(x,t)\,d\mu\) converges uniformly in $t$, then
\[
\int_\alpha^\omega\int_{\mathcal M}f\,d\mu\,dt
=\int_{\mathcal M}\int_\alpha^\omega f\,dt\,d\mu.
\]
\end{lemma}
\begin{proof}
\sketch Uniform convergence makes the tails outside large balls uniformly small.  Apply
Fubini on each ball and pass to the limit, with the tail estimate controlling
the error.
\end{proof}

\begin{remark}[Uniform tail criterion]
\label{rem:chow-iii-ch23-fubini-tail}
\source{chow-et-al-ricci-flow-part-iii:page-0282}
\dcref{ch23:23.39}
\group{Supplementary analysis}

The uniform convergence hypothesis means that for every $\varepsilon>0$ there
is an $R$ such that
\(\left|\int_{\mathcal M\setminus B(O,R)}f(x,t)\,d\mu\right|<\varepsilon\)
for all $t\in[\alpha,\omega]$.
\end{remark}

\begin{lemma}[Differentiation under the integral sign]
\label{lem:chow-iii-ch23-differentiation-under-integral}
\source{chow-et-al-ricci-flow-part-iii:page-0283,chow-et-al-ricci-flow-part-iii:page-0284}
\dcref{ch23:23.40}
\group{Supplementary analysis}

If $f\in C^1(\mathcal M\times[\alpha,\omega])$ and both
\(\int f\,d\mu\) and $\int\partial_tf\,d\mu$ converge uniformly in $t$, then
\[
\frac d{dt}\int_{\mathcal M}f(x,t)\,d\mu(x)
=\int_{\mathcal M}\partial_tf(x,t)\,d\mu(x).
\]
\end{lemma}
\begin{proof}
\sketch Apply Lemma~\ref{lem:chow-iii-ch23-fubini-noncompact} to $\partial_tf$ and use
the fundamental theorem of calculus on each spatial point; uniform convergence
permits passage to the spatial integral.
\end{proof}

\begin{example}[Problem 23.36 (Asymptotic entropy)]
\label{ex:chow-iii-ch23-entropy-problem}
\dcref{ch23:23.36}
\group{Asymptotics}
\source{chow-et-al-ricci-flow-part-iii:page-0280}
Determine the asymptotics as $t\to0$ of $-\int_{\mathcal M}fH\,d\mu$ for
$H=(4\pi t)^{-n/2}e^{-f}$.
\end{example}
